\documentclass[../../../../SoftwareRequirementsSpecification.tex]{subfiles}
\begin{document}
% Richiede le macro definite in src/macros/useCaseMacros.tex
% (incluse nel preambolo di SoftwareRequirementsSpecification.tex)

\ucstart
\begin{xltabular}{\textwidth}{|l|X|}
  \hline
  \ucid{A-02} \\ \hline
  \ucname{Richiesta Appuntamento} \\ \hline
  \ucactor{Atleta} \\ \hline
  \ucdescription{L'atleta richiede un appuntamento con il PT.
    L'atleta indica la palestra, la data, l'ora e la motivazione.} \\
    \hline
  \ucprecond{L'atleta deve essere autenticato.} \\ \hline
  \ucpostcond{Il sistema crea un appuntamento in attesa
    di conferma dal PT. Il PT riceve un'email di notifica.} \\ \hline
  \ucmainflow{
    \item L'atleta apre la pagina dei propri appuntamenti.
    \item L'atleta preme il pulsante "Richiedi Appuntamento".
    \item Il sistema mostra il form di richiesta appuntamento, con
      l'elenco delle palestre a cui l'atleta è iscritto (vedi Nota 1).
    \item L'atleta inserisce i dati necessari: la palestra in cui
      vuole allenarsi, la data e l'ora dell'appuntamento e la
      motivazione.
    \item L'atleta preme il pulsante "Invia Richiesta".
    \item Il sistema verifica i dati nel database.
    \item Il sistema salva nel database la richiesta di appuntamento.
    \item Il sistema invia al PT un'email per informarlo
      della nuova richiesta.
  } \\ \hline
  \ucaltflow{
    \textbf{6a. Richiesta in attesa già presente} \newline
    - Il sistema mostra un messaggio di errore ("C'è già una richiesta
    di appuntamento in attesa di risposta."). \newline
    - Il flusso riprende dal passo 1. \newline
  } \\ \hline
  \ucnote{
    \textbf{Nota 1:} se l'atleta è iscritto a una sola palestra, il
    sistema la preseleziona e l'atleta può comunque confermarla o
    cambiarla, se iscritto a più di una (vedi caso d'uso A-06).
  } \\ \hline
\end{xltabular}
\end{document}
